\begin{thebibliography}{99}
\raggedright
\addcontentsline{toc}{section}{References}

\bibitem{ProveIt}
V.~Reshetnikov and ProveIt contributors,
\emph{ProveIt: Unknot recognition research and incoming reports}.
Repository snapshot \baseline, inspected 9 October 2026.
\href{https://github.com/VladimirReshetnikov/ProveIt/tree/3a90fb34146c915328ab8eac6250cc2514f74ed0/Topology/UnknotRecognition}{Pinned unknot-recognition source tree}.
Incoming directory:
\href{https://github.com/VladimirReshetnikov/ProveIt/tree/3a90fb34146c915328ab8eac6250cc2514f74ed0/docs/incoming}{Pinned incoming directory}.

\bibitem{priorKernel}
ProveIt research continuation,
\emph{Quadrilateral Support Kernels and Certified Normal-Disc Discovery},
9 October 2026. Incoming archive
\code{unknot_quadrilateral_support_kernel_20261009.zip} at the pinned snapshot.
\href{https://github.com/VladimirReshetnikov/ProveIt/blob/3a90fb34146c915328ab8eac6250cc2514f74ed0/docs/incoming/unknot_quadrilateral_support_kernel_20261009.zip}{Pinned source archive}.
The contraction, mixed-minor estimate, canonical ray lifting and earlier layered support obstruction are dependencies of this continuation.

\bibitem{HLP}
J.~Hass, J.~C. Lagarias and N.~Pippenger,
\emph{The computational complexity of knot and link problems},
Journal of the ACM \textbf{46} (1999), no.~2, 185--211.
\url{https://arxiv.org/abs/math/9807016}.

\bibitem{Tollefson}
J.~L. Tollefson,
\emph{Normal surface Q-theory},
Pacific Journal of Mathematics \textbf{183} (1998), no.~2, 359--374.
\url{https://msp.org/pjm/1998/183-2/pjm-v183-n2-p07-s.pdf}.

\bibitem{BurtonQ}
B.~A. Burton,
\emph{Converting between quadrilateral and standard solution sets in normal surface theory},
Algebraic \& Geometric Topology \textbf{9} (2009), 2121--2174.
\url{https://arxiv.org/abs/0901.2629}.

\bibitem{Suh}
C.~Suh,
\emph{The unknotting problem and normal surface Q-theory},
arXiv:1009.1500, first submitted 8 September 2010.
\url{https://arxiv.org/abs/1009.1500}.

\bibitem{BO}
B.~A. Burton and M.~\"Ozlen,
\emph{A fast branching algorithm for unknot recognition with experimental polynomial-time behaviour},
arXiv:1211.1079v3, 9 October 2014.
\url{https://arxiv.org/pdf/1211.1079}.
Numbering in this report refers to version~3: Algorithm~9, Lemmas~5--10 and Theorems~11--12.

\bibitem{AHT}
I.~Agol, J.~Hass and W.~P. Thurston,
\emph{The computational complexity of knot genus and spanning area},
arXiv:math/0205057v2, 10 July 2004.
\url{https://arxiv.org/abs/math/0205057}.

\bibitem{BurtonHe}
B.~A. Burton and A.~He,
\emph{On the hardness of finding normal surfaces},
Journal of Applied and Computational Topology \textbf{5} (2021), 583--619.
\url{https://arxiv.org/abs/1912.09051v3}.

\bibitem{Khachiyan}
L.~G. Khachiyan,
\emph{A polynomial algorithm in linear programming},
Doklady Akademii Nauk SSSR \textbf{244} (1979), no.~5, 1093--1096.
\url{https://www.mathnet.ru/eng/dan42319}.

\bibitem{Bland}
R.~G. Bland,
\emph{New finite pivoting rules for the simplex method},
Mathematics of Operations Research \textbf{2} (1977), no.~2, 103--107.
\url{https://doi.org/10.1287/moor.2.2.103}.

\bibitem{LackenbyTalk}
M.~Lackenby,
\emph{Unknot recognition in quasi-polynomial time},
Oxford seminar slides, March 2021, 34-page handout.
\url{https://people.maths.ox.ac.uk/lackenby/quasipolynomial-talk-oxford-compressed.pdf}.
The displayed announced bound is $2^{O((\log n)^3)}$.

\bibitem{Lackenby2026}
M.~Lackenby,
\emph{Incompressible surfaces, hierarchies and unknot recognition},
arXiv:2607.23350v1, 25 July 2026.
\url{https://arxiv.org/abs/2607.23350}.
Section~9 distinguishes iteration bounds from an implementation running-time estimate; this report does not import a quasi-polynomial runtime theorem from that paper.

\end{thebibliography}
